
\subsection*{A5: Dynamic Group Chat 
} \label{sec:app:groupchat}

\newcommand{\alice}{\text{Alice}\xspace}
\newcommand{\bob}{\text{Bob}\xspace}

\newcommand{\userproxy}{\text{User Proxy}\xspace}


\libName provides native support for a \emph{dynamic group chat} communication pattern, in which participating agents share the same context and converse with the others in a dynamic manner instead of following a pre-defined order. 
Dynamic group chat relies on ongoing conversations to guide the flow of interaction among agents. 
These make dynamic group chat ideal for situations where collaboration without strict communication order is beneficial.  
In \libName, the \groupchatmanager class serves as the conductor of conversation among agents and repeats the following three steps: dynamically selecting a speaker, collecting responses from the selected speaker, and broadcasting the message (Figure~\ref{fig:app_summary}-A5). For the dynamic speaker-selection component, we use a role-play style prompt. 
Through a pilot study on 12 manually crafted complex tasks, we observed that compared to a prompt that is purely based on the task, utilizing a role-play prompt often leads to more effective consideration of both conversation context and role alignment during the problem-solving and speaker-selection process. Consequently, this leads to a higher success rate and fewer LLM calls. We include detailed results in Appendix~\ref{appendix:app:groupchat}.